 \subsection{解法の変更}

    先に挙げたチャートによってではうまく解消されない時は一度最初からやり直してみるとよい.   
    この際に今までと同じ考え方ではなく,なるべく別のアプローチができないか検討すると解法を見つけやすい.

    また,それ以外の場合でも次のどれかが起こったら,今の方法を一度止め,解法の再検討を行うと良い.
    \begin{enumerate}
      \item 対数を取ったが,対数の中身が一致しない
      \item 特性方程式の解が存在しない.（虚数解を除く）
      \item 置換後に新しい種類の項が増え続ける.
    \end{enumerate}
    このようなときには$a_2$から$a_5$ぐらいまでの値を求めてみると法則性に気づき,帰納法で証明できる場合がある.
    

    それでも見つけられない時は特殊関数型漸化式であるか,突飛な発想を求められるものである.
    一旦ご飯でも食べて改めて取り組もう.
    ご飯の時間でないのなら,今までの数学の知識を総動員して似たものを探そう.
    
    この手順を身につけると,個別のパターンを暗記するだけでなく,
    見慣れない漸化式にも最初の一手を選べるようになる.